\documentclass[11pt]{article}
\usepackage[margin=1in,headheight=14pt]{geometry}
\usepackage[T1]{fontenc}
\usepackage{xcolor}
\usepackage{listings}
\usepackage{fancyhdr}
\usepackage[hidelinks]{hyperref}
\usepackage{parskip}

\definecolor{codebg}{RGB}{245,245,245}
\definecolor{codeframe}{RGB}{200,200,200}
\definecolor{katzblue}{RGB}{30,80,140}

\lstset{
  backgroundcolor=\color{codebg},
  basicstyle=\ttfamily\small,
  breaklines=true,
  frame=single,
  rulecolor=\color{codeframe},
  xleftmargin=10pt,
  xrightmargin=10pt,
  framesep=6pt,
  columns=fullflexible,
  keepspaces=true
}

\pagestyle{fancy}
\fancyhf{}
\fancyhead[L]{KatzLab Ureteroscope --- Mint-Only Guide}
\fancyhead[R]{\thepage}
\renewcommand{\headrulewidth}{0.4pt}

\title{\textbf{\color{katzblue} KatzLab Ureteroscope}\\[4pt]
Doing Everything on the Mint Box Alone}
\author{No Mac, no second machine, no Foxglove --- one machine, start to finish}
\date{\today}

\begin{document}
\maketitle
\thispagestyle{fancy}

\section*{Why this guide is different from the main runbook}
\texttt{docs/RUNBOOK.pdf} assumes two machines: the Mint box running the
rig and Isaac Sim, and a Mac watching it remotely over Foxglove. That
split exists only because Isaac Sim cannot run on macOS at all --- it
needs an NVIDIA GPU and Linux/Windows.

If you are sitting at the Mint box itself, none of that split is needed.
Everything --- the bridge, RViz, Isaac Sim, the real rig CAD, training ---
runs and is viewed locally, on one machine, with no WebSocket bridge, no
second app, and no network setup. This guide is that simpler path.

\tableofcontents
\newpage

%============================================================
\section{Install ROS 2}
%============================================================

\subsection{Check the Mint/Ubuntu base}
The install script refuses to run on anything other than a Mint 22.x
(Ubuntu 24.04 ``noble'') base, on purpose:

\begin{lstlisting}
cat /etc/upstream-release/lsb-release
\end{lstlisting}

You want \texttt{DISTRIB\_CODENAME=noble}. If it says something else,
stop and ask before continuing.

\subsection{Run the setup script}
From the repository root:

\begin{lstlisting}
cd katzlab-ureteroscope
./ros2/setup-mint.sh
\end{lstlisting}

Installs ROS 2 Jazzy Desktop (RViz2 included), \texttt{colcon}, and
\texttt{foxglove-bridge} (unused in this guide, but harmless to have). A
few GB, takes a while, safe to re-run if it fails partway.

\subsection{Open a new terminal and verify}
A terminal that was already open will not see the newly-installed ROS 2.
Open a fresh one:

\begin{lstlisting}
ros2 --version
which colcon
rviz2 --help | head -1
\end{lstlisting}

All three should print something, not ``command not found.''

%============================================================
\section{Sanity-check what needs nothing else first}
%============================================================

These run right now, before ROS 2 or Isaac Sim are even touched. If
either fails, fix it before going further --- nothing downstream will
work:

\begin{lstlisting}
cd ros2/katzlab_bridge/test && python3 -m pytest -q
cd ../../isaac/training && python3 -m pytest test_reach_math.py -q
\end{lstlisting}

Expected: \texttt{6 passed, 1 skipped} and \texttt{10 passed}.

%============================================================
\section{Build the ROS 2 package}
%============================================================

\begin{lstlisting}
cd ros2
colcon build --packages-select katzlab_bridge
source install/setup.bash
\end{lstlisting}

Do this in every new terminal you open for the rest of this guide (or add
it to \texttt{\textasciitilde/.bashrc} once you are confident it works).

%============================================================
\section{Bridge + RViz, all local}
%============================================================

One command starts the bridge node and RViz together, no Foxglove, no
WebSocket, nothing but this one machine:

\begin{lstlisting}
ros2 launch katzlab_bridge sim.launch.py dry_run:=true
\end{lstlisting}

\texttt{dry\_run:=true} means no Teensy/hardware needed yet --- this
exercises the ROS 2 wiring only. Drop the flag once hardware is attached.

An RViz window opens directly on this machine's own screen, showing the
placeholder stick-figure model. Leave this terminal running.

\subsection{Send a test motion command}
From a second terminal:

\begin{lstlisting}
source ~/katzlab-ureteroscope/ros2/install/setup.bash

ros2 topic pub /ureteroscope/cmd_velocity geometry_msgs/msg/Twist \
  "{linear: {x: 0.001}, angular: {z: 0.0, y: 0.0}}"
\end{lstlisting}

The model in RViz should move. Other useful checks from this same
terminal:

\begin{lstlisting}
ros2 topic echo /ureteroscope/joint_states
ros2 service call /ureteroscope/home std_srvs/srv/Trigger
ros2 service call /ureteroscope/estop std_srvs/srv/Trigger
\end{lstlisting}

%============================================================
\section{Isaac Sim: the real rig CAD}
%============================================================

This is the part that is genuinely local-only no matter what --- Isaac
Sim cannot be watched remotely from a Mac at all, so being at this
machine is the only way to see it, full stop.

\subsection{Look at it before driving it}
\begin{lstlisting}
source /opt/ros/jazzy/setup.bash
~/isaacsim/python.sh ~/katzlab-ureteroscope/ros2/isaac/preview_scene.py
\end{lstlisting}

A window opens with the real rig CAD (linear stage, rotation bearing, the
actual Wolf RIWO ureteroscope mesh) and the real 8-stone benchtop target
grid. Check:

\begin{enumerate}
  \item Does the scope's length look visually right relative to the
        table/linear stage? (There is a known open question about this
        --- see \texttt{ros2/isaac/assets/ASSET-NOTES.md}.)
  \item Do the 7 flexible-tip segments bend as a plausible continuous
        curve when driven through their range?
  \item Do the 8 stone markers sit somewhere sensible near the tip?
\end{enumerate}

If something looks visually wrong here, stop and note exactly what, before
trusting anything built on top of it.

\subsection{Drive it live from the bridge}
With Section~4's \texttt{sim.launch.py} still running in its own
terminal:

\begin{lstlisting}
~/isaacsim/python.sh ~/katzlab-ureteroscope/ros2/isaac/load_rig.py
\end{lstlisting}

Re-send the test command from Section~4.1. Both the RViz window
(placeholder model) and the Isaac Sim window (real CAD) should move
together, within a step or two of each other.

%============================================================
\section{First training milestone: get the tip to the stone}
%============================================================

\subsection{Smoke-test}
\begin{lstlisting}
~/isaacsim/python.sh \
  ~/katzlab-ureteroscope/ros2/isaac/training/reach_task.py
\end{lstlisting}

Random actions, no training --- confirms the environment runs end to end.

\subsection{Install training dependencies}
\begin{lstlisting}
~/isaacsim/python.sh -m pip install gymnasium numpy stable-baselines3
\end{lstlisting}

\subsection{Train}
\begin{lstlisting}
~/isaacsim/python.sh \
  ~/katzlab-ureteroscope/ros2/isaac/training/train_reach.py \
  --total-timesteps 200000
\end{lstlisting}

State-based (tip/target \emph{positions}, not camera images) on purpose
--- simpler problem first, and the prerequisite for a vision-based
version later. Do not expect reliable stone-reaching from one run; watch
for mean episode reward trending up and mean episode length trending down
as the first sign anything is working.

%============================================================
\section{Everything running at once --- a reference layout}
%============================================================

Four terminals (or four panes in one terminal multiplexer), all on the
Mint box:

\begin{center}
\begin{tabular}{|l|l|}
\hline
\textbf{Terminal} & \textbf{Running} \\
\hline
1 & \texttt{ros2 launch katzlab\_bridge sim.launch.py dry\_run:=true} \\
2 & \texttt{\textasciitilde/isaacsim/python.sh .../load\_rig.py} \\
3 & free, for \texttt{ros2 topic pub} / \texttt{service call} commands \\
4 & free, for training once 1--3 are confirmed working \\
\hline
\end{tabular}
\end{center}

No Mac, no Foxglove, no second machine, no network configuration --- one
box, four terminals, everything visible on its own screen.

%============================================================
\section{What to report back if something breaks}
%============================================================

For each step: which numbered section, the exact command, and the full
error output (not a paraphrase).

Section~5 (Isaac Sim) is the one most likely to need a fix on first
contact with a real install --- none of that code has been run against
one yet. The two most likely failure points (the URDF importer
extension's exact module path, and Replicator/RL API field names) are
both named in \texttt{ros2/isaac/README.md}.

\end{document}
